\section{真正的新缩放轴：同样的数据，怎样学得更好}

这堂课把预训练从“模型有多大、token 有多少”的问题，推进成“\textbf{模型以什么顺序、什么目标、什么反馈学习这些 token}”的问题。讲者串起三项工作：two-phase pretraining 决定数据混合与顺序，front-loading reasoning 决定推理数据何时进入训练，RLP（Reinforcement as a Pretraining Objective）则把探索式思考直接放进预训练目标。三者都在改变学习过程，而不是简单增加语料。

\subsection{四个组件不是四张独立清单}

课程先给出一个 SOTA LLM 的四组件配方。Smart Data 负责质量、覆盖与去重；Smart Architecture 负责表示与计算结构；Smart Algorithms 决定优化目标、课程和训练阶段；Smart Collaboration 则让 pretraining、post-training、research 与 engineering 不再各自优化局部指标。读图时要注意：本讲重点是第三项，但任何算法收益都依赖另外三项提供可用数据、稳定模型与可执行系统。

\lecturefigure{slide-002.jpg}{构建 SOTA LLM 的四个耦合组件}{Stanford 官方录像，00:01:21--00:02:09}

\teachervoice{00:01:21--00:02:09，讲者把 smart collaboration 与 data、architecture、algorithm 并列。这个口头提醒很重要：预训练 recipe 不是一篇论文里的单一超参数，而是多个团队共享数据定义、算力预算和评测口径后的系统结果。}

\begin{importantbox}{本讲的总命题}
当模型规模和数据规模继续上升时，新的能力不只来自“更多”，还来自\textbf{更有结构的学习过程}：先看什么数据、何时看推理数据、是否允许模型在预测前探索，以及奖励是否能把有效探索与无效自言自语区分开。
\end{importantbox}

\subsection{从 smart data 转向 smart algorithms}

讲者先用一页概括其数据工作，包括数学对话、Common Crawl 数学抽取、多样化 synthetic-data generation 与跨域 reasoning data。这里不是要逐一讲数据集，而是说明后面的算法实验并非在真空中进行：数据质量标签、领域覆盖与推理样本定义，都是训练策略可以操作的输入变量。最终展开状态还提示这些资源有公开版本，便于外部复核。

\lecturefigure{slide-003.jpg}{Smart Data 工作概览：数学、网页抽取、多样化合成与跨域推理}{Stanford 官方录像，00:02:09--00:03:15}

接下来，课程明确切换到三个算法问题。第一项问“怎样把既有数据潜力用满”，第二项问“推理数据是否应该更早出现”，第三项问“能否在预训练时让模型通过探索获得奖励”。三项工作逐层深入：先改变采样分布，再改变阶段分配，最后改变训练目标。若把它们混成一个 recipe，就无法判断收益究竟来自数据、顺序还是 objective。

\lecturefigure{slide-004.jpg}{本讲三项算法工作：two-phase、front-loading 与 RLP}{Stanford 官方录像，00:03:15--00:03:37}

\begin{knowledgebox}{术语消化：三个干预层级}
\begin{tabularx}{\textwidth}{L{0.22\textwidth}X X}
\toprule
干预 & 改变什么 & 保持什么相对可比 \\
\midrule
Optimal blend & 各数据源权重与重复次数 & 总 token budget 与模型配置 \\
Two-phase curriculum & 数据在训练进度中的出现顺序 & 总体数据池与目标函数 \\
Front-loading reasoning & 推理数据进入训练流程的阶段 & 固定总推理数据预算与后训练流程 \\
RLP & 预测前是否采样 thought、怎样计算 reward & 真实下一 token 与通用预训练流 \\
\bottomrule
\end{tabularx}
\end{knowledgebox}

\subsection{Pascal、Volta、Ampere、Hopper：控制变量的教学隐喻}

四个学习者拥有相同数据，却采用不同策略。Pascal 不做 curriculum、不使用 reasoning-rich data、也不通过思考探索；Volta 只做 curriculum；Ampere 再加入早期 reasoning data；Hopper 最后加入 learning through thinking。最终矩阵不是儿童教育理论，而是一套实验索引：后续每张结果图都在比较某一列从叉号变成勾号后发生什么。

\lecturefigure{slide-005.jpg}{四个学习者对应的 curriculum、front-loading 与 learning-through-thinking 策略}{Stanford 官方录像，00:03:37--00:06:27}

\teachervoice{00:04:15--00:05:18，讲者用“随机读一篇历史、一道高阶数学、再读基础数学”的例子解释无 curriculum，又用提前修 AP 课程解释 front-loading。隐喻的价值在于区分“拥有数据”和“在正确阶段使用数据”。}

\begin{warningbox}{隐喻不是因果证据}
Pascal 到 Hopper 的人物故事只帮助记忆实验条件。真实模型并不理解“先打基础”的教育学语义；是否有效必须回到受控数据预算、训练 compute、benchmark 与置信区间。后文所有数字都只描述论文中的具体设置。
\end{warningbox}

\subsection{本章小结}

本章建立了全课比较坐标：smart data 提供可操作的数据源，smart algorithms 决定权重、顺序、阶段和奖励；Pascal 到 Hopper 则把三次干预组织成递进实验。下面先处理最保守的一步——不改模型与损失，只重新设计数据 mixture 与 curriculum。

